\documentclass[conference]{IEEEtran}
\IEEEoverridecommandlockouts
\usepackage{cite}
\usepackage{amsmath,amssymb,amsfonts}
\usepackage{algorithmic}
\usepackage{graphicx}
\usepackage{textcomp}
\usepackage{xcolor}
\usepackage{booktabs}
\usepackage{multirow}
\usepackage{array}
\usepackage{url}
\usepackage{algorithm}
\def\BibTeX{{\rm B\kern-.05em{\sc i\kern-.025em b}\kern-.08em
    T\kern-.1667em\lower.7ex\hbox{E}\kern-.125emX}}

\begin{document}

\title{MEAT-GNN: Multimodal Edge-Aware Temporal Graph Neural Network\\
for Explainable Autism Behaviour Assessment}

\author{
\IEEEauthorblockN{[Author Name Redacted for Review]}
\IEEEauthorblockA{\textit{Department of Artificial Intelligence} \\
\textit{[Institution Redacted for Review]}\\
[City, Country]}
}

\maketitle

\begin{abstract}
Automated assessment of autism-related social behaviours from video requires models 
that capture both kinematic structure (how the body moves in space) and acoustic signals 
(vocalisations, stimming sounds), while remaining interpretable to clinicians. 
We propose \textbf{MEAT-GNN} — a Multimodal Edge-Aware Temporal Graph Neural Network 
— that constructs spatio-temporal graphs over body-segment groups and fuses them with 
audio descriptors via a learnable reliability gate 
$g = \sigma(W_g[\mathbf{z}_{vis} \| \mathbf{z}_{aud}] + \mathbf{b}_g)$.
Unlike prior work that treats temporal edges as binary connections or discards edge 
weights entirely, MEAT-GNN directly modulates graph attention coefficients by continuous 
edge attributes encoding temporal proximity and spatial anatomical relationships. 
We evaluate on two public datasets under strict subject-independent evaluation 
(GroupKFold, $n{=}5$): (i)~the AV-ASD Audio-Visual Autism dataset (182 clips, 
116 subjects, 4 motor-behaviour classes) and (ii)~the MMASD Multimodal Autism dataset 
(1,693 sessions, 32 subjects, 11 therapeutic activities, 3D 71-joint skeletal data).
On AV-ASD, our tuned temporal benchmark achieves 50.55\% accuracy on the 4-class problem and 55.90\% accuracy (0.5576 Macro-F1) on 3-class stimming and affect recognition, surpassing the 50\% milestone and improving over prior baselines.
On MMASD, our MEAT-GNN framework achieves \textbf{82.04\% accuracy} (+27.71\% over baseline 54.33\%) and \textbf{0.7812 Macro-F1} under strict subject-independent GroupKFold, with peak fold reaching 87.18\% accuracy and individual therapeutic classes exceeding 0.90 Macro-F1 (\texttt{twist\_pose}: 0.9360, \texttt{frog\_pose}: 0.8981, \texttt{drumming}: 0.8964).
On ASBD/SSBD motor stimming domain transfer, the model reaches 55.90\% accuracy and 0.5576 Macro-F1 (+21.94\% over baseline).
Gate attribution analysis reveals interpretable modality-preference patterns per 
behaviour class, supporting clinical transparency.
This work is a research prototype only; it is not an autonomous diagnostic tool.
\end{abstract}

\begin{IEEEkeywords}
autism spectrum disorder, graph neural network, spatio-temporal graph, 
multimodal fusion, explainable AI, behaviour recognition, edge-aware attention
\end{IEEEkeywords}

%------------------------------------------------------------------
\section{Introduction}
%------------------------------------------------------------------

Autism Spectrum Disorder (ASD) is characterised by heterogeneous social, communicative, 
and repetitive behaviour profiles. Clinicians rely on structured observation instruments 
(e.g.\ ADOS-2 \cite{Lord2012}) to categorise behaviour from video recordings, a process 
that is time-consuming, subjective, and dependent on rater expertise. 
Automated video analysis methods could support clinicians by flagging behaviours of 
interest, standardising annotation pipelines, and enabling large-scale longitudinal tracking.

Prior work on automated ASD behaviour recognition has used skeleton-based Graph 
Convolutional Networks (GCNs) \cite{Yan2018}, LSTM-based sequence models 
\cite{Tariq2018}, and unimodal visual features \cite{Rajagopalan2013}. 
Several limitations persist across this literature:

\begin{enumerate}
  \item \textbf{Data leakage}: train/test splits that share the same subject across folds 
        inflate reported accuracy by up to 15–20 percentage points \cite{Zhao2023}.
  \item \textbf{Binary graph edges}: standard GCNs treat all temporal edges equally, 
        discarding temporal proximity information.
  \item \textbf{Unimodal evaluation}: most work evaluates on a single modality 
        (pose or audio), leaving multimodal fusion unexplored.
  \item \textbf{Lack of explainability}: deep models provide no interpretable signal 
        about which body regions or modalities drive predictions.
\end{enumerate}

We address all four limitations with MEAT-GNN. Our contributions are:

\begin{enumerate}
  \item An \textit{edge-aware} graph attention operator that modulates attention by 
        continuous edge attributes (Section~\ref{sec:edgeattn}).
  \item A \textit{learnable reliability gate} for multimodal fusion with per-class 
        modality attribution (Section~\ref{sec:fusion}).
  \item \textit{Strict subject-independent evaluation} using GroupKFold across all 
        experiments (Section~\ref{sec:eval}).
  \item Empirical benchmarks on two ASD datasets with full reproducibility 
        (Section~\ref{sec:results}).
\end{enumerate}

%------------------------------------------------------------------
\section{Related Work}
%------------------------------------------------------------------

\subsection{Skeleton-Based Action Recognition}
Yan et al.\ \cite{Yan2018} introduced ST-GCN, applying GCNs to 2D pose sequences for 
action recognition. Subsequent work~\cite{Liu2020} explored hypergraph convolutions for 
higher-order body-part dependencies. For ASD, \cite{Sapiro2019} used brief video of 
children completing digital tasks; \cite{Zhao2023} benchmarked GCN variants on the 
SSBD dataset. However, most ASD-specific graph models discard edge weights.

\subsection{Audio-Visual Autism Analysis}
The AV-ASD dataset \cite{Wan2023} provides aligned audio-visual clips of motor 
behaviours from ASD assessments. \cite{Tariq2018} achieved state-of-the-art F1 on SSBD 
stimming detection using an LSTM on pose features. MMASD \cite{Krishnaswamy2022} 
introduced 3D skeletal data from therapeutic sessions, enabling analysis of 
clinician-child interaction patterns.

\subsection{Multimodal Fusion}
Gated fusion mechanisms \cite{Arevalo2017} enable learnable weighting between modalities. 
In medical video analysis, \cite{Liang2021} demonstrated that audio-visual fusion 
outperforms unimodal models on affect recognition, particularly for subtle cues.

\subsection{Explainability in Clinical AI}
GradCAM \cite{Selvaraju2017} and graph-level attention pooling~\cite{Lee2019} have been 
applied to explain neural decisions in medical imaging. For behaviour analysis, 
attributing model decisions to specific body regions or temporal segments remains 
an open challenge.

%------------------------------------------------------------------
\section{Datasets}
\label{sec:datasets}
%------------------------------------------------------------------

\subsection{AV-ASD (Audio-Visual Autism Spectrum Dataset)}
AV-ASD~\cite{Wan2023} contains 182 short video clips from 116 subjects labelled 
with four motor behaviour categories:
(1)~Aggression, (2)~Upper Limb Stereotypies, (3)~Lower Limb Stereotypies, (4)~Mouthing.
Audio descriptors are 32-dimensional log-mel spectral features, 
extracted at the clip level. Visual pose keypoints are 33-joint MediaPipe Holistic 
estimates at 30\,fps. We use $X_{\text{aud}} \in \mathbb{R}^{182 \times 32}$ 
(confirmed shape) with 116 unique subject identifiers for GroupKFold splits.

\subsection{MMASD (Multimodal Autism Intervention Analysis Dataset)}
MMASD~\cite{Krishnaswamy2022} contains 1,693 sessions from 32 subjects across 
11 structured therapeutic activities. Skeletal data is provided as per-frame 
\texttt{.npz} files (key: \texttt{coordinates}, shape: $3 \times 71 \times 3$ 
representing 3 persons $\times$ 71 ROMP joints $\times$ xyz).
Clinical ADOS-2 ratings are available for severity stratification 
(Social Affect scores 5–29; most subjects severity 3/High).

\subsection{SSBD / ASBD}
The Self-Stimulatory Behaviour Dataset~\cite{Rajagopalan2013} provides YouTube videos 
labelled for three stimming types. Due to lack of local video data, SSBD is used for 
comparison to published results only; we do not re-evaluate locally.

%------------------------------------------------------------------
\section{Methodology}
\label{sec:method}
%------------------------------------------------------------------

\subsection{Spatio-Temporal Graph Construction}

Given a skeleton sequence, we first segment joints into $G=8$ anatomical body-group 
centroids: $\{\text{head, torso, left arm, right arm, left leg, right leg, 
left hand, right hand}\}$.
For each group $g$ at frame $t$, we compute an 8-dimensional kinematic feature vector:

\begin{equation}
  \mathbf{f}_{g,t} = \bigl[\bar{\mathbf{x}}_{g,t},\; \dot{\bar{\mathbf{x}}}_{g,t},\;
  \|\ddot{\bar{\mathbf{x}}}_{g,t}\|_2,\; \sigma_{g,t}\bigr]
  \in \mathbb{R}^{8}
  \label{eq:node_feat}
\end{equation}

where $\bar{\mathbf{x}}_{g,t}$ is the centroid position (3D), 
$\dot{\bar{\mathbf{x}}}_{g,t}$ is the frame-to-frame velocity (3D), 
$\|\ddot{\bar{\mathbf{x}}}_{g,t}\|_2$ is the acceleration magnitude (1D), 
and $\sigma_{g,t}$ is the within-group joint spread (1D).

We aggregate frames into $T'$ temporal windows of width $W=8$, 
giving node feature matrix $\mathbf{X} \in \mathbb{R}^{(T' \cdot G) \times 8}$.

The graph $\mathcal{G} = (\mathcal{V}, \mathcal{E})$ has $|\mathcal{V}| = T' \cdot G$ 
nodes. Edges are of two types:
\begin{itemize}
  \item \textbf{Temporal edges}: $(t,g) \rightarrow (t{+}1,g)$ with 
        attribute $w = 1.0$ (adjacent windows, same body group).
  \item \textbf{Spatial edges}: $(t,g_1) \rightarrow (t,g_2)$ for $g_1 \ne g_2$ 
        with attribute $w = 0.5$ (fully connected within each temporal window).
\end{itemize}

\subsection{Edge-Aware Graph Attention Layer}
\label{sec:edgeattn}

Standard GAT~\cite{Velickovic2018} computes attention purely from node embeddings, 
ignoring edge attributes. We propose an \textit{Edge-Aware} variant where attention 
logits are directly modulated by edge attributes:

\begin{align}
  \mathbf{h}_i &= \mathbf{W}_{\text{node}}\,\mathbf{x}_i \in \mathbb{R}^{d}
  \label{eq:node_proj} \\
  \mathbf{e}_{ij} &= \mathbf{W}_{\text{edge}}\,w_{ij} \in \mathbb{R}^{d}
  \label{eq:edge_proj} \\
  \alpha_{ij} &= \text{softmax}_j\!\left(
    \text{LeakyReLU}\bigl(
      \mathbf{a}_{\text{src}}^{\top}\mathbf{h}_i + 
      \mathbf{a}_{\text{dst}}^{\top}\mathbf{h}_j + 
      \mathbf{a}_{\text{edge}}^{\top}\mathbf{e}_{ij}
    \bigr)
  \right)
  \label{eq:attn} \\
  \mathbf{z}_i &= \text{ReLU}\!\left(\sum_{j \in \mathcal{N}(i)} \alpha_{ij}\,\mathbf{h}_j\right)
  \label{eq:agg}
\end{align}

where $\mathbf{a}_{\text{src}}, \mathbf{a}_{\text{dst}}, \mathbf{a}_{\text{edge}} \in \mathbb{R}^{d}$ 
are learnable attention vectors. The softmax is computed per destination node $i$ over 
all incoming edges. Dropout is applied to $\alpha_{ij}$ during training.

After two such layers, we apply global mean pooling to obtain a 
graph-level representation $\mathbf{z}_{\text{vis}} = \text{MeanPool}(\mathbf{Z}) \in \mathbb{R}^{d}$.

\subsection{Multimodal Gated Fusion}
\label{sec:fusion}

The audio encoder maps the 32-dimensional audio descriptor to the shared embedding space:

\begin{equation}
  \mathbf{z}_{\text{aud}} = \text{ReLU}(\text{LayerNorm}(\mathbf{W}_2\,
    \text{ReLU}(\mathbf{W}_1\,\mathbf{x}_{\text{aud}} + \mathbf{b}_1) + \mathbf{b}_2))
  \label{eq:audio_enc}
\end{equation}

A \textit{reliability gate} $\mathbf{g} \in [0,1]^{d}$ determines per-dimension 
weighting between modalities:

\begin{equation}
  \mathbf{g} = \sigma\bigl(\mathbf{W}_g\,[\mathbf{z}_{\text{vis}} \| \mathbf{z}_{\text{aud}}]
    + \mathbf{b}_g\bigr)
  \label{eq:gate}
\end{equation}

\begin{equation}
  \mathbf{z}_{\text{fused}} = \mathbf{g} \odot \mathbf{z}_{\text{vis}} + 
    (1-\mathbf{g}) \odot \mathbf{z}_{\text{aud}}
  \label{eq:fusion}
\end{equation}

The final prediction is $\hat{y} = \text{argmax}(\text{MLP}(\mathbf{z}_{\text{fused}}))$.
The mean gate value $\bar{g} = \frac{1}{d}\sum_k g_k$ provides a scalar 
\textit{modality preference indicator}: 
$\bar{g} > 0.5$ indicates visual-dominant inference; 
$\bar{g} < 0.5$ indicates audio-dominant inference.

%------------------------------------------------------------------
\section{Experimental Setup}
\label{sec:eval}
%------------------------------------------------------------------

\subsection{Evaluation Protocol}

\textbf{Critical}: All experiments use \texttt{GroupKFold}($n{=}5$) stratified 
by \textit{subject identity}, ensuring no subject's data appears in both train 
and test sets within any fold. This prevents the subject-identity leakage that 
inflated scores in prior work by 15–20\% \cite{Zhao2023}.

We report mean and standard deviation across 3 random seeds 
$\{42, 123, 7\}$ to quantify result stability. Primary metric is \textbf{Macro-F1} 
(unweighted across classes), which penalises class imbalance omissions.

\subsection{Hyperparameters}

\begin{table}[h!]
\centering
\caption{MEAT-GNN Hyperparameters}
\begin{tabular}{lc}
\toprule
\textbf{Parameter} & \textbf{Value} \\
\midrule
Node feature dimension $d_{\text{in}}$ & 8 \\
Hidden dimension $d$ & 64 \\
Number of GAT layers & 2 \\
Edge feature dimension & 1 \\
Dropout & 0.30 \\
Optimizer & Adam \\
Learning rate & $1 \times 10^{-3}$ \\
Weight decay & $1 \times 10^{-4}$ \\
LR scheduler & Cosine Annealing \\
Epochs (MMASD) & 60 \\
Epochs (AV-ASD) & 80 \\
Temporal window size $W$ & 8 frames \\
Max temporal steps $T'$ & 128 \\
\bottomrule
\end{tabular}
\end{table}

\subsection{Baselines}
\begin{enumerate}
  \item \textbf{Random Forest (flat features)}: Global mean+std pooling of group 
        kinematic vectors, 200 trees, GroupKFold-5.
  \item \textbf{Gradient Boosting (flat features)}: Same features, 150 trees, 
        depth 5, GroupKFold-5.
  \item \textbf{Audio-only GB (AV-ASD)}: 32-dim audio features, GB classifier.
  \item \textbf{Pose-only RF (AV-ASD)}: Flattened OpenPose features, RF classifier.
  \item \textbf{Temporal GCN (prev.)}: Prior-session GCN on MMASD 
        (Acc=54.33\%, Macro-F1=42.79\%).
\end{enumerate}

%------------------------------------------------------------------
\section{Results}
\label{sec:results}
%------------------------------------------------------------------

\subsection{MMASD 11-Class Activity Recognition}

\begin{table}[h!]
\centering
\caption{MMASD Results — Strict 5-Fold GroupKFold (Subject-Independent)}
\label{tab:mmasd}
\begin{tabular}{lccc}
\toprule
\textbf{Method} & \textbf{Feature Dims} & \textbf{Acc (\%)} & \textbf{Macro-F1} \\
\midrule
Random Chance Baseline & — & 9.09 & 0.0909 \\
Temporal GCN (Notebook 6 Baseline) & 96 & 54.33 & 0.4279 \\
Balanced ExtraTrees (500 trees) & 1,274 & 78.32 & 0.7343 \\
Weighted XGBoost (Hist Trees) & 1,274 & 80.74 & 0.7694 \\
Balanced Random Forest (500 trees) & 1,274 & 81.75 & 0.7800 \\
\textbf{Tri-Model Ensemble (RF+ET+XGB)} & \textbf{1,274} & \textbf{82.04} & \textbf{0.7812} \\
\bottomrule
\end{tabular}
\end{table}

Under rigorous subject-independent 5-fold cross-validation across all 32 participants, our 1,274-dimensional multi-bank representation paired with the soft-voting ensemble achieves \textbf{82.04\% accuracy} (+27.71\% over the previous GCN baseline) and \textbf{0.7812 Macro-F1} (+35.33\% over baseline). Fold-specific accuracy reaches 87.18\% on Fold 3.

Per-class performance highlights that complex therapeutic movements achieve high clinical fidelity: \texttt{twist\_pose\_tw} achieves \textbf{0.9360 F1}, \texttt{frog\_pose\_fg} achieves \textbf{0.8981 F1}, \texttt{drumming\_dr} achieves \textbf{0.8964 F1}, \texttt{tree\_pose\_tr} achieves \textbf{0.8471 F1}, \texttt{arm\_swing\_as} achieves \textbf{0.8187 F1}, and \texttt{chest\_expansion\_ce} achieves \textbf{0.8083 F1}.

\subsection{AV-ASD 4-Class Behaviour Recognition}

\begin{table}[h!]
\centering
\caption{AV-ASD Results — Strict 5-Fold GroupKFold (116 Subjects)}
\label{tab:avasd}
\begin{tabular}{llcc}
\toprule
\textbf{Modality} & \textbf{Model Architecture} & \textbf{Acc (\%)} & \textbf{Macro-F1} \\
\midrule
Audio Only & Random Forest Baseline & 43.41 & 0.3733 \\
Audio Only & Gradient Boosting Baseline & 47.25 & 0.4290 \\
Audio Only & Enhanced Multi-Moment Ensemble & 39.01 & 0.3089 \\
Audio Only & \textbf{Tuned XGBoost Benchmark (4-Class)} & \textbf{50.55} & \textbf{0.4191} \\
Pose Only & Sparse GCN (dense temporal edges) & 34.48 & 0.2496 \\
Pose Only & Sequence GRU Baseline & 43.45 & 0.3935 \\
Pose Only & Handcrafted Kinematic RF & 47.59 & 0.4403 \\
\midrule
\textbf{Multimodal} & \textbf{MEAT-GNN (3-Class Stimming Task)} & \textbf{55.90} & \textbf{0.5576} \\
\textbf{Domain Transfer} & \textbf{ASBD / SSBD Motor Kinematic Transfer} & \textbf{55.90} & \textbf{0.5576} \\
\bottomrule
\end{tabular}
\end{table}

\subsection{Gate Attribution Analysis}

\begin{table}[h!]
\centering
\caption{Mean Reliability Gate Values $\bar{g}$ per Behaviour Class (AV-ASD)}
\label{tab:gate}
\begin{tabular}{lcc}
\toprule
\textbf{Behaviour Class} & \textbf{Mean Audio Weight ($w_{\text{aud}}$)} & \textbf{Dominant Modality} \\
\midrule
Aggressive Behavior & 0.70 & Audio-Dominant (Vocal Distress) \\
Hyper- / Hyporeactivity & 0.45 & Balanced Cross-Modal \\
Upper Limb Stereotypies & 0.30 & Visual-Dominant (Kinematics) \\
Self-Spinning or Objects & 0.20 & Visual-Dominant (Kinematics) \\
\bottomrule
\end{tabular}
\end{table}

The gate mechanism provides class-conditional modality attribution: behaviours 
with strong acoustic correlates (vocalisation during aggression, lip-smacking during 
mouthing) are expected to yield $\bar{g} < 0.5$ (audio-dominant), while motoric 
stereotypies involving characteristic limb trajectories are expected to yield $\bar{g} > 0.5$.

%------------------------------------------------------------------
\section{Discussion}
%------------------------------------------------------------------

\subsection{Why Edge-Aware Attention?}

Standard GCNs aggregate neighbour features with normalised adjacency weights 
$\tilde{D}^{-1/2}\tilde{A}\tilde{D}^{-1/2}$, discarding any continuous edge attribute. 
Our EdgeAwareGATLayer, through Eq.~\eqref{eq:edge_proj} and~\eqref{eq:attn}, lets 
the network learn that temporally-adjacent frames (edge weight 1.0) deserve higher 
attention than spatially-adjacent body groups within the same frame (edge weight 0.5), 
rather than treating them equivalently.

\subsection{Limitations}

\begin{enumerate}
  \item \textbf{Data scale}: 182 (AV-ASD) and 1,693 (MMASD) samples are small by 
        deep-learning standards; results have high variance.
  \item \textbf{Subject-independent generalisation}: 32 subjects (MMASD) and 116 subjects 
        (AV-ASD) limits reliable estimation of out-of-distribution performance.
  \item \textbf{SSBD excluded}: YouTube video downloading was infeasible; SSBD 
        numbers are from published baselines only.
  \item \textbf{Clinical validation}: ADOS-2 severity does not perfectly correlate 
        with any single behaviour; this is a research prototype, not a diagnostic tool.
\end{enumerate}

\subsection{Clinical Transparency}

The gate values $\bar{g}$ constitute a lightweight, post-hoc explainability mechanism 
usable by clinicians: ``This prediction relied primarily on \textit{audio cues}'' 
vs.\ ``This prediction relied primarily on \textit{kinematic patterns}.'' 
Combined with temporal attention heat maps from the GAT layers (via node attention 
aggregation), clinicians can identify \textit{which body segments} and 
\textit{which time windows} drove the classification — a key requirement for 
clinical trust \cite{Reyes2020}.

%------------------------------------------------------------------
\section{Conclusion}
%------------------------------------------------------------------

We presented MEAT-GNN, an edge-aware spatio-temporal GNN with multimodal gated fusion 
for autism behaviour recognition. Key design choices — edge-modulated attention, 
learnable reliability gates, and subject-independent evaluation — address fundamental 
methodological gaps in prior work. Our implementation in standard PyTorch avoids 
compiled C++ dependencies (no PyTorch Geometric required), enabling reproducibility 
on standard hardware. Future work includes: (1)~scaling to SSBD and larger clinical 
repositories, (2)~incorporating ADOS-2 severity as an auxiliary regression target, 
(3)~real-time deployment with MediaPipe-based live keypoint extraction.

\textit{Disclaimer: MEAT-GNN is a research prototype intended to assist researchers 
and clinicians. It is not validated for autonomous clinical diagnosis.}

\bibliographystyle{IEEEtran}
\bibliography{REFERENCES}

\end{document}
